\documentclass[11pt]{article}
\usepackage[margin=1in]{geometry}
\usepackage{amsmath,amsfonts,amssymb,amsthm,enumerate,bbm}
\usepackage[]{graphicx}
\usepackage{color,subfigure}
\definecolor{scarlet}{rgb}{0.81,0,0}
\usepackage{multicol}
\usepackage{float}
\usepackage[all]{xypic}
\usepackage[colorlinks=true,citecolor=scarlet,linkcolor=scarlet]{hyperref}
\usepackage{colonequals}

\usepackage{fancyhdr, lastpage}
\pagestyle{fancy}
\fancyfoot[C]{{\thepage} of \pageref{LastPage}}



\DeclareMathOperator{\mSpec}{mSpec}
\DeclareMathOperator{\Spec}{Spec}
\DeclareMathOperator{\Ass}{Ass}
\DeclareMathOperator{\Supp}{Supp}
\DeclareMathOperator{\height}{height}
\DeclareMathOperator{\Hom}{Hom}
\DeclareMathOperator{\ann}{ann}
\DeclareMathOperator{\End}{End}
\DeclareMathOperator{\coker}{coker}
%\DeclareMathOperator{\ker}{ker}
\DeclareMathOperator{\rank}{rank}
\DeclareMathOperator{\im}{im}
\DeclareMathOperator{\M}{M}
\DeclareMathOperator{\Tor}{Tor}
\DeclareMathOperator{\id}{id}
\DeclareMathOperator{\ch}{char}
\DeclareMathOperator{\Aut}{Aut}

%\DeclareMathOperator{\dim}{dim}

\DeclareMathOperator{\lcm}{lcm}

\def\ra{\rightarrow}
\newcommand{\m}{\mathfrak{m}}
\newcommand{\C}{\mathbb{C}}
\newcommand{\Q}{\mathbb{Q}}
\newcommand{\Z}{\mathbb{Z}}
\newcommand{\ZZ}{\mathbb{Z}}
\newcommand{\R}{\mathbb{R}}
\newcommand{\N}{\mathbb{N}}
\newcommand{\ov}[1]{\overline{#1}}
\newcommand{\norm}{\trianglelefteq}

\def\ov#1{\overline{#1}}


\title{}
\date{\vspace{-0.5in}}

\makeatletter
\g@addto@macro\@floatboxreset\centering
\makeatother

\theoremstyle{definition}
\newtheorem{problem}{Problem}


\begin{document}

\thispagestyle{fancy}
\pagestyle{fancy}
\rhead{UNL $\mid$ Spring 2026}
\lhead{Introduction to Modern Algebra II}

\vspace{3em}

\begin{center}
	{\LARGE Problem Set 10 \\}
	Due Thursday, April 9\end{center}

\

\noindent
{\bf Instructions:}
You are encouraged to work together on these problems, but each student should hand in their own final draft, written in a way that indicates their individual understanding of the solutions. Never submit something for grading that you do not completely understand. You cannot use any resources besides me, your classmates, and our course notes.


I will post the .tex code for these problems for you to use if you wish to type your homework. If you prefer not to type, please  {\em write neatly}. As a matter of good proof writing style, please use complete sentences and correct grammar. You may use any result stated or proven in class or in a homework problem, provided you reference it appropriately by either stating the result or stating its name (e.g. the definition of ring or Lagrange's Theorem). Please do not refer to theorems by their number in the course notes, as that can change.



\begin{problem}
Let $E$ be the field extension of $\Q$ obtained by adjoining to $\Q$ all four complex roots of the polynomial $x^4 + 5$.
That is, $E = \Q(\alpha_1, \alpha_2, \alpha_3, \alpha_4)$ where
$$\alpha_1 = e^{\pi i/4}\sqrt[4]{5}, \quad
\alpha_2 = e^{3\pi i/4}\sqrt[4]{5}, \quad
\alpha_3 = e^{5\pi i/4}\sqrt[4]{5}, \quad
\alpha_4 = e^{7\pi i/4}\sqrt[4]{5}.$$

\begin{enumerate}[(a)]
\item Prove that there exists\footnote{Note that $\alpha_1 + \alpha_4$ is a real number; find it explicitly.} 
 a field extension $\Q \subseteq F$ such that $F \subseteq E$, $F \subseteq \R$, and $[F :\Q]=4$.


\item Determine $[E : \Q]$ with justification.
\end{enumerate}
\end{problem}


\begin{problem}
Let $F$ be a field and $f,g \in F[x]$ be nonzero polynomials. Show that $\gcd(f,g) = 1$ in $F[x]$ if and only if $f$ and $g$ have no common roots in an algebraic closure $\overline{F}$ of $F$.
\end{problem}



\begin{problem}
	Let $F \subseteq L$ be a field extension and let $S \subseteq L$ be an arbitrary subset of $L$ whose elements are all algebraic over $F$. Show\footnote{You can use without proof the analogue of Lemma 15.1.12 for larger generating sets: if $F\subseteq L$ is a field extension and $S\subseteq L$, then $F(S) = \{ \frac{f(s_1,\dots,s_t)}{g(s_1,\dots,s_t)} \ | \ s_1,\dots,s_t\in S, \ f(x_1,\dots,x_t), g(x_1,\dots,x_t)\in F[x_1,\dots,x_t], \ g(s_1,\dots,s_t)\neq 0\}$.}
	 that $F \subseteq F(S)$ is algebraic.
\end{problem}


\begin{problem}
	Let $L$ be the splitting field of $f = x^5 - 11 \in \Q[x]$.
	\begin{enumerate}[(a)]
		\item Find the degree\footnote{Consider the chains $\Q \subseteq \Q(\sqrt[5]{11}) \subseteq L$ and $\Q \subseteq \Q(\xi) \subseteq L$.} of $[L : \Q]$.
		\item Let $F = \Q(\xi)$, where $\xi = e^{\frac{2 \pi i}{5}}$ is a primitive fifth root of unity. Show that $f$ is irreducible over~$F$.
	\end{enumerate}
\end{problem}


\end{document}